\subsection{Related Works}

Our work builds on the intersection of strategic machine learning, sender-receiver games, and voluntary disclosure in predictive modeling. Early studies in strategic machine learning, such as \cite{hardt2016strategic}, formalized settings where agents manipulate features to influence outcomes, showing that robust models can mitigate gaming in linear settings. This was extended in \cite{ghalme2021strategic} to dark (unknown) environments, where agents' responses are not fully observable, leading to equilibrium analyses under uncertainty.

Sender-receiver games provide a foundational framework for our analysis. \cite{blume2002learning} examined learning dynamics in these games using experimental data, highlighting how repeated interactions evolve into efficient signaling. Evolutionary perspectives, as in \cite{skyrms2009evolution}, explore multi-agent signaling, while \cite{blume1998learning} compared stimulus-response and belief-based learning frameworks.

Voluntary disclosure literature addresses selective revelation. \cite{mcdonough2023voluntary} analyzed investor disclosures, predicting good-news bias, which parallels our heterogeneous preferences. Empirical studies, such as \cite{lu2024determinants}, used XGBoost and XAI for predicting disclosure patterns, emphasizing transparency in strategic contexts.

Stackelberg equilibria are central to our game-theoretic approach. \cite{jin2022incentive} characterized equilibria in strategic regression, designing incentives to prevent cheating. \cite{sundaram2023pac} generalized PAC-learning for heterogeneous agents. In adversarial prediction, \cite{bruckner2011stackelberg} framed problems as Stackelberg games, solving for optimal decision rules.

While these works advance strategic aspects in regression \cite{shavit2020causal}, they often assume known environments or linear models without noisy reporting channels. Our contributions provide formal sufficiency and necessity conditions under noisy channels and heterogeneous preferences, filling a gap in generalized predictive settings.
